% GENERATED FILE. Edit canonical lessons, then run npm run generate:books.
% canonical-source-hash: fnv1a64:0abdbc51d875d8c9
% canonical-lessons: UR-C223-chalak, UR-C223-sharif, UR-C223-badtamiz, UR-C223-tamizdar, UR-C223-hasin

\chapter{Cunning, Decent, Rude, Polite, Beautiful}
\label{ch:ur-cunning-decent-rude-polite-beautiful}
\hlchaptermodality{223}

\begin{chapteropening}
\textbf{By the end of this chapter:} \emph{I can say cunning, decent, rude, polite and beautiful.}

Complete the last lesson of chapter 223: I can say cunning, decent, rude, polite and beautiful.
\end{chapteropening}

\section[\texorpdfstring{\ur{چالاک} (chālāk)}{چالاک (chālāk)}]{\ur{چالاک} (chālāk) — cunning}

\label{lesson:UR-C223-chalak}



\begin{quote}
Before the new one: say the Urdu for fake, then the Urdu for truthful.
\end{quote}

\subsection*{You'll want to know: \ur{چالاک}}
\textbf{\ur{چالاک}} — \emph{chālāk} — \textquotedblleft{}cunning\textquotedblright{}.

\begin{cousinweb}
\textbf{\ur{چالاک}} (\emph{chālāk}) is a loan, usually traced to Persian \emph{chālāk}, \textquotedblleft{}quick, clever\textquotedblright{}.
\end{cousinweb}

\begin{grammarlens}[title={What you've built}]
One more word for this chapter.
\end{grammarlens}

\subsection*{Your turn}
\begin{itemize}
\raggedright
  \item \emph{Say it:} \emph{chālāk}
  \item \emph{Say it:} \emph{chālāk}, once more
  \item \emph{Say it:} say \emph{chālāk}
  \item \emph{Recall:} say the Urdu for fake, then the Urdu for truthful, then say \emph{chālāk} again
\end{itemize}

\begin{tcolorbox}[breakable,colback=teal!4,colframe=teal!35!black,title={Before you move on}]
What does \ur{چالاک} mean? (\textquotedblleft{}Cunning\textquotedblright{}.) Say it once more.
\end{tcolorbox}
\section[\texorpdfstring{\ur{شریف} (sharīf)}{شریف (sharīf)}]{\ur{شریف} (sharīf) — decent, well-mannered}

\label{lesson:UR-C223-sharif}



\begin{quote}
Before the new one: say the Urdu for truthful, then the Urdu for cunning.
\end{quote}

\subsection*{You'll want to know: \ur{شریف}}
\textbf{\ur{شریف}} — \emph{sharīf} — \textquotedblleft{}decent, well-mannered\textquotedblright{}.

\begin{cousinweb}
\textbf{\ur{شریف}} (\emph{sharīf}) is a loan from Arabic \emph{sharīf}, \textquotedblleft{}noble, of high rank\textquotedblright{}, which reached Urdu through Persian. Its Arabic root is \emph{sh-r-f} (Lane).
\end{cousinweb}

\begin{grammarlens}[title={What you've built}]
One more word for this chapter.
\end{grammarlens}

\subsection*{Your turn}
\begin{itemize}
\raggedright
  \item \emph{Say it:} \emph{sharīf}
  \item \emph{Say it:} \emph{sharīf}, once more
  \item \emph{Say it:} say \emph{sharīf}
  \item \emph{Recall:} say the Urdu for truthful, then the Urdu for cunning, then say \emph{sharīf} again
\end{itemize}

\begin{tcolorbox}[breakable,colback=teal!4,colframe=teal!35!black,title={Before you move on}]
What does \ur{شریف} mean? (\textquotedblleft{}Decent, well-mannered\textquotedblright{}.) Say it once more.
\end{tcolorbox}
\section[\texorpdfstring{\ur{بدتمیز} (badtamīz)}{بدتمیز (badtamīz)}]{\ur{بدتمیز} (badtamīz) — rude}

\label{lesson:UR-C223-badtamiz}



\begin{quote}
Before the new one: say the Urdu for cunning, then the Urdu for decent.
\end{quote}

\subsection*{You'll want to know: \ur{بدتمیز}}
\textbf{\ur{بدتمیز}} — \emph{badtamīz} — \textquotedblleft{}rude\textquotedblright{}.

\begin{cousinweb}
\textbf{\ur{بدتمیز}} (\emph{badtamīz}) is \textquotedblleft{}bad in manners\textquotedblright{}. It joins Persian \emph{bad}, \textquotedblleft{}bad\textquotedblright{}, and \emph{tamīz}, \textquotedblleft{}manners\textquotedblright{}, Arabic \emph{tamyīz}, \textquotedblleft{}discernment, telling apart\textquotedblright{} (Lane, root \emph{m-y-z}).
\end{cousinweb}

\begin{grammarlens}[title={What you've built}]
One more word for this chapter.
\end{grammarlens}

\subsection*{Your turn}
\begin{itemize}
\raggedright
  \item \emph{Say it:} \emph{badtamīz}
  \item \emph{Say it:} \emph{badtamīz}, once more
  \item \emph{Say it:} say \emph{badtamīz}
  \item \emph{Recall:} say the Urdu for cunning, then the Urdu for decent, then say \emph{badtamīz} again
\end{itemize}

\begin{tcolorbox}[breakable,colback=teal!4,colframe=teal!35!black,title={Before you move on}]
What does \ur{بدتمیز} mean? (\textquotedblleft{}Rude\textquotedblright{}.) Say it once more.
\end{tcolorbox}
\section[\texorpdfstring{\ur{تمیزدار} (tamīzdār)}{تمیزدار (tamīzdār)}]{\ur{تمیزدار} (tamīzdār) — polite}

\label{lesson:UR-C223-tamizdar}



\begin{quote}
Before the new one: say the Urdu for decent, then the Urdu for rude.
\end{quote}

\subsection*{You'll want to know: \ur{تمیزدار}}
\textbf{\ur{تمیزدار}} — \emph{tamīzdār} — \textquotedblleft{}polite\textquotedblright{}.

\begin{cousinweb}
\textbf{\ur{تمیزدار}} (\emph{tamīzdār}) is \textquotedblleft{}holding manners\textquotedblright{}. \emph{tamīz}, \textquotedblleft{}manners\textquotedblright{}, is Arabic \emph{tamyīz}, \textquotedblleft{}discernment, telling apart\textquotedblright{} (Lane, root \emph{m-y-z}), and \emph{-dār}, \textquotedblleft{}holder\textquotedblright{}, is Persian.
\end{cousinweb}

\begin{grammarlens}[title={What you've built}]
One more word for this chapter.
\end{grammarlens}

\subsection*{Your turn}
\begin{itemize}
\raggedright
  \item \emph{Say it:} \emph{tamīzdār}
  \item \emph{Say it:} \emph{tamīzdār}, once more
  \item \emph{Say it:} say \emph{tamīzdār}
  \item \emph{Recall:} say the Urdu for decent, then the Urdu for rude, then say \emph{tamīzdār} again
\end{itemize}

\begin{tcolorbox}[breakable,colback=teal!4,colframe=teal!35!black,title={Before you move on}]
What does \ur{تمیزدار} mean? (\textquotedblleft{}Polite\textquotedblright{}.) Say it once more.
\end{tcolorbox}
\section[\texorpdfstring{\ur{حسین} (hasīn)}{حسین (hasīn)}]{\ur{حسین} (hasīn) — beautiful}

\label{lesson:UR-C223-hasin}



\begin{quote}
Before the new one: say the Urdu for rude, then the Urdu for polite.
\end{quote}

\subsection*{You'll want to know: \ur{حسین}}
\textbf{\ur{حسین}} — \emph{hasīn} — \textquotedblleft{}beautiful\textquotedblright{}.

\begin{cousinweb}
\textbf{\ur{حسین}} (\emph{hasīn}) is a loan from Arabic \emph{ḥasīn}, \textquotedblleft{}beautiful, comely\textquotedblright{}, which reached Urdu through Persian. Its Arabic root is \emph{ḥ-s-n} (Lane). It is built on the same root as \emph{ḥusn}, \textquotedblleft{}beauty\textquotedblright{}.
\end{cousinweb}

\begin{grammarlens}[title={What you've built}]
One more word for this chapter.
\end{grammarlens}

\subsection*{Your turn}
\begin{itemize}
\raggedright
  \item \emph{Say it:} \emph{hasīn}
  \item \emph{Say it:} \emph{hasīn}, once more
  \item \emph{Say it:} say \emph{hasīn}
  \item \emph{Recall:} say the Urdu for rude, then the Urdu for polite, then say \emph{hasīn} again
\end{itemize}

\begin{tcolorbox}[breakable,colback=teal!4,colframe=teal!35!black,title={Before you move on}]
What does \ur{حسین} mean? (\textquotedblleft{}Beautiful\textquotedblright{}.) Say it once more.
\end{tcolorbox}
